\section{Trabalhos relacionados}
\subsection{Avaliação de alucinação e capacidade}
POPE~\cite{pope} avalia presença de objetos por perguntas binárias e escolhe objetos ausentes aleatórios, frequentes ou coocorrentes com os presentes. RePOPE~\cite{repope} revisa seus rótulos e remove perguntas ambíguas, permitindo examinar o efeito de ruído de anotação nas mesmas respostas. CHAIR~\cite{chair} identifica objetos mencionados sem suporte nas anotações em descrições livres. Sua versão por descrição penaliza a ocorrência de qualquer objeto alucinado; a versão por instância considera a proporção de menções alucinadas. Como a métrica não fixa imagens, instrução nem extensão, adotamos o protocolo registrado a partir de OPERA~\cite{opera}. MME~\cite{mme} complementa esses instrumentos com subtarefas de percepção e cognição. COCO~\cite{coco} fornece as imagens e anotações usadas nas tarefas de presença e descrição e nos ajustes prescritos pelas receitas.

\subsection{Intervenções nas representações}
VTI~\cite{vti} constrói direções visuais e textuais e intervém nas camadas do codificador e do decodificador. REVIS~\cite{revis} adiciona um vetor visual em uma camada, condicionado a um limiar de risco calibrado. VISTA~\cite{vista} deriva um vetor por imagem da diferença entre execuções com e sem imagem e combina logits intermediários. Nullu~\cite{nullu} projeta ativações fora de um subespaço associado a descrições alucinadas. Os dois últimos estão implementados no projeto, mas não têm avaliação principal com pesos reais nos resultados aqui analisados.

V-ITI~\cite{viti} e RUDDER~\cite{rudder} aparecem nos documentos como comparações publicadas de mitigação. Seus protocolos e modelos de referência diferem, e nenhum dos dois avalia VTI ou REVIS na comparação registrada. Assim, não usamos os números entre artigos como uma classificação comum nem atribuímos a esses trabalhos mecanismos que a documentação disponível não descreve.

\subsection{Leitura e intervenção nos estados internos}
A \textit{logit lens}, leitura de estados intermediários pela cabeça de saída, indica quando a resposta se torna acessível ao vocabulário final. Uma \textit{probe}, ou sonda linear, verifica se o rótulo já pode ser decodificado antes dessa etapa. São leituras distintas: a primeira usa a cabeça existente; a segunda aprende uma leitura supervisionada.

\textit{Attention knockout}, bloqueio de rotas de atenção, e \textit{activation patching}, substituição de ativações por estados de uma execução de referência, examinam dependências por intervenção. Complementamos esses procedimentos com a atenção ponderada pela norma do que cada cabeça escreve, conforme a abordagem de Kobayashi et al.~\cite{kobayashi}. A ponderação distingue peso de atenção de magnitude transmitida, mas continua sendo uma leitura correlacional. A combinação dessas leituras permite distinguir informação decodificável, dependência de uma rota e associação com regiões visuais.
